\section{Error analysis across modalities}
\label{app:error_analysis_all}

To complement the final test scores and the learning curves, we also inspect the
label-wise error structure of all models across all five modalities.
For each model and modality, the corresponding plots report the missed-positive
error rate for each functional group, defined as the proportion of false
negatives relative to the number of positive test examples for that label.
This view is useful because aggregate micro-F1 can remain comparatively strong
even when several functional groups are predicted poorly.
All error-analysis plots are reported on the held-out test set rather than on
the validation set, since the validation split is already used for model
selection and early stopping.

\paragraph{IR.}
IR yields the most balanced error profiles overall.
This is consistent with the main quantitative comparison, where the adapted CNN
 is strongest and MPS is the closest tensor-network competitor.
Across all four models, the same groups repeatedly appear among the most
difficult labels, most notably thial, azo compound, hydrazone, phosphine, and
sulfoxide.
The key point is therefore not that one model exhibits a unique IR-specific
failure mode, but that the residual IR errors are concentrated in a small set of
chemically difficult labels whose spectral signatures overlap strongly with
related groups.

\begin{figure}[H]
    \centering
    \begin{subfigure}[t]{0.48\linewidth}
        \centering
        \includegraphics[width=\linewidth]{Error_Analysis/CNN_IR_test_error_analysis_detailed.png}
        \caption{CNN}
    \end{subfigure}
    \hfill
    \begin{subfigure}[t]{0.48\linewidth}
        \centering
        \includegraphics[width=\linewidth]{Error_Analysis/MPS_IR_test_error_analysis_detailed.png}
        \caption{MPS}
    \end{subfigure}

    \vspace{0.6em}

    \begin{subfigure}[t]{0.48\linewidth}
        \centering
        \includegraphics[width=\linewidth]{Error_Analysis/TTN_IR_test_error_analysis_detailed.png}
        \caption{TTN}
    \end{subfigure}
    \hfill
    \begin{subfigure}[t]{0.48\linewidth}
        \centering
        \includegraphics[width=\linewidth]{Error_Analysis/XGB_IR_test_error_analysis_detailed.png}
        \caption{XGBoost}
    \end{subfigure}
    \caption{Detailed test-set error analysis for IR across all evaluated models.}
    \label{fig:error_analysis_ir_all}
\end{figure}

\paragraph{$^{1}$H-NMR.}
For $^{1}$H-NMR, all models show substantially higher label-wise error than on
IR, with a visibly larger mass of labels in the high-error regime.
The adapted CNN remains the most balanced model overall, while TTN and XGBoost
show broader tails of severe per-label errors and MPS remains weaker on many
minority labels.
At the same time, the difficult functional groups are again highly consistent
across models, especially sulfonic acid, thial, isothiocyanate, isocyanate, and
thiol.
This suggests that the main challenge is not merely optimization instability,
but the limited separability of several functional groups under the present
$^{1}$H-NMR representation.

\begin{figure}[H]
    \centering
    \begin{subfigure}[t]{0.48\linewidth}
        \centering
        \includegraphics[width=\linewidth]{Error_Analysis/CNN_H-NMR_test_error_analysis_detailed.png}
        \caption{CNN}
    \end{subfigure}
    \hfill
    \begin{subfigure}[t]{0.48\linewidth}
        \centering
        \includegraphics[width=\linewidth]{Error_Analysis/MPS_H-NMR_test_error_analysis_detailed.png}
        \caption{MPS}
    \end{subfigure}

    \vspace{0.6em}

    \begin{subfigure}[t]{0.48\linewidth}
        \centering
        \includegraphics[width=\linewidth]{Error_Analysis/TTN_H-NMR_test_error_analysis_detailed.png}
        \caption{TTN}
    \end{subfigure}
    \hfill
    \begin{subfigure}[t]{0.48\linewidth}
        \centering
        \includegraphics[width=\linewidth]{Error_Analysis/XGB_H_NMR_test_error_analysis_detailed.png}
        \caption{XGBoost}
    \end{subfigure}
    \caption{Detailed test-set error analysis for $^{1}$H-NMR across all evaluated models.}
    \label{fig:error_analysis_hnmr_all}
\end{figure}

\paragraph{$^{13}$C-NMR.}
The $^{13}$C-NMR error analysis mirrors the main results in the body of the
paper.
XGBoost shows the most favorable overall label-wise error profile, while both
tensor-network models remain clearly more robust than the adapted CNN, whose
performance collapses for many minority labels.
MPS and TTN reduce the number of near-complete failures relative to the CNN, but
still retain a substantial set of difficult functional groups.
The most problematic labels again include sulfonic acid, thial,
isothiocyanate, azo compound, and phosphine, indicating that the difficulty of
$^{13}$C-NMR is concentrated in a recurring subset of rare and chemically
specific targets.

\begin{figure}[H]
    \centering
    \begin{subfigure}[t]{0.48\linewidth}
        \centering
        \includegraphics[width=\linewidth]{Error_Analysis/CNN_C-NMR_test_error_analysis_detailed.png}
        \caption{CNN}
    \end{subfigure}
    \hfill
    \begin{subfigure}[t]{0.48\linewidth}
        \centering
        \includegraphics[width=\linewidth]{Error_Analysis/MPS_C-NMR_test_error_analysis_detailed.png}
        \caption{MPS}
    \end{subfigure}

    \vspace{0.6em}

    \begin{subfigure}[t]{0.48\linewidth}
        \centering
        \includegraphics[width=\linewidth]{Error_Analysis/TTN_C-NMR_test_error_analysis_detailed.png}
        \caption{TTN}
    \end{subfigure}
    \hfill
    \begin{subfigure}[t]{0.48\linewidth}
        \centering
        \includegraphics[width=\linewidth]{Error_Analysis/XGB_C_NMR_test_error_analysis_detailed.png}
        \caption{XGBoost}
    \end{subfigure}
    \caption{Detailed test-set error analysis for $^{13}$C-NMR across all evaluated models.}
    \label{fig:error_analysis_cnmr_all}
\end{figure}

\paragraph{MS/MS$^{+}$.}
Positive-ion MS/MS is markedly more difficult than IR and both NMR modalities.
Among the four models, MPS exhibits the least collapsed label-wise error
distribution, whereas CNN, TTN, and XGBoost place a large fraction of
functional groups close to complete missed-positive failure.
This observation is consistent with the learning curves, where no model obtains
uniformly strong performance across both common and rare labels.
Under the current one-dimensional representation, MPS preserves some
minority-label signal, but the modality remains challenging for all evaluated
architectures.

\begin{figure}[H]
    \centering
    \begin{subfigure}[t]{0.48\linewidth}
        \centering
        \includegraphics[width=\linewidth]{Error_Analysis/CNN_MSMS+_test_error_analysis_detailed.png}
        \caption{CNN}
    \end{subfigure}
    \hfill
    \begin{subfigure}[t]{0.48\linewidth}
        \centering
        \includegraphics[width=\linewidth]{Error_Analysis/MPS_MSMS+_test_error_analysis_detailed.png}
        \caption{MPS}
    \end{subfigure}

    \vspace{0.6em}

    \begin{subfigure}[t]{0.48\linewidth}
        \centering
        \includegraphics[width=\linewidth]{Error_Analysis/TTN_MSMS+_test_error_analysis_detailed.png}
        \caption{TTN}
    \end{subfigure}
    \hfill
    \begin{subfigure}[t]{0.48\linewidth}
        \centering
        \includegraphics[width=\linewidth]{Error_Analysis/XGB_MSMS+_test_error_analysis_detailed.png}
        \caption{XGBoost}
    \end{subfigure}
    \caption{Detailed test-set error analysis for MS/MS$^{+}$ across all evaluated models.}
    \label{fig:error_analysis_msms_pos_all}
\end{figure}

\paragraph{MS/MS$^{-}$.}
Negative-ion MS/MS shows essentially the same global pattern as positive-ion
MS/MS.
MPS again yields the most graded error distribution, while the remaining three
models are dominated by severe per-label failures.
The persistence of this structure across both ionization modes suggests that the
main bottleneck is not a single architecture choice, but the difficulty of
representing sparse fragmentation spectra in the current fixed-length
one-dimensional setup.
For this reason, improvements on MS/MS are likely to depend at least as much on
representation design as on the choice of classifier.

\begin{figure}[H]
    \centering
    \begin{subfigure}[t]{0.48\linewidth}
        \centering
        \includegraphics[width=\linewidth]{Error_Analysis/CNN_MSMS-_test_error_analysis_detailed.png}
        \caption{CNN}
    \end{subfigure}
    \hfill
    \begin{subfigure}[t]{0.48\linewidth}
        \centering
        \includegraphics[width=\linewidth]{Error_Analysis/MPS_MSMS-_test_error_analysis_detailed.png}
        \caption{MPS}
    \end{subfigure}

    \vspace{0.6em}

    \begin{subfigure}[t]{0.48\linewidth}
        \centering
        \includegraphics[width=\linewidth]{Error_Analysis/TTN_MSMS-_test_error_analysis_detailed.png}
        \caption{TTN}
    \end{subfigure}
    \hfill
    \begin{subfigure}[t]{0.48\linewidth}
        \centering
        \includegraphics[width=\linewidth]{Error_Analysis/XGB_MSMS-_test_error_analysis_detailed.png}
        \caption{XGBoost}
    \end{subfigure}
    \caption{Detailed test-set error analysis for MS/MS$^{-}$ across all evaluated models.}
    \label{fig:error_analysis_msms_neg_all}
\end{figure}

Overall, the error-analysis plots reinforce three conclusions from the main
paper.
First, IR is the most favorable modality, with relatively balanced error
profiles and shared residual failure modes across models.
Second, both NMR modalities remain substantially harder, especially for rare and
chemically specific functional groups.
Third, the MS/MS modalities are the most difficult setting in this benchmark,
and their weakness is broad enough across models that improving the input
representation appears at least as important as changing the classifier.
Across all modalities, labels such as sulfonic acid, thial, azo compound,
hydrazone, isothiocyanate, and phosphine repeatedly appear among the most
difficult groups, indicating that the dominant sources of error are tied to the
chemical structure of the benchmark rather than to a single model family.
